% GENERATED FILE. Edit canonical lessons, then run npm run generate:books.
% canonical-source-hash: fnv1a64:90a519ba7da38994
% canonical-lessons: RU-C407-obgonyat, RU-C407-perestraivatsya, RU-C407-svorachivat, RU-C407-razvorachivatsya, RU-C407-obochina

\chapter{Overtaking and turning}
\label{ch:ru-overtaking-and-turning}
\hlchaptermodality{407}

\begin{chapteropening}
\textbf{By the end of this chapter:} \emph{I can describe moves on the road: overtaking, changing lanes, turning off, turning round and the hard shoulder.}

Complete the last lesson of chapter 407: I can describe moves on the road: overtaking, changing lanes, turning off, turning round and the hard shoulder.
\end{chapteropening}

\section[\texorpdfstring{\ru{обгонять} (obgonyát)}{обгонять (obgonyát)}]{\ru{обгонять} (obgonyát) — to overtake}

\label{lesson:RU-C407-obgonyat}



\begin{quote}
Before the new one: say the Russian for a seat, then the Russian for a windscreen wiper.
\end{quote}

\subsection*{You'll want to know: \ru{обгонять}}
\textbf{\ru{обгонять}} (\emph{obgonyát}) — \textquotedblleft{}to overtake\textquotedblright{}.

\begin{quote}
\ru{На} \ru{этой} \ru{дороге} \ru{опасно} \ru{обгонять}.
\end{quote}

\begin{quote}
\emph{Na étoy doróge opásno obgonyát.}
\end{quote}

\begin{quote}
It is dangerous to overtake on this road.
\end{quote}

Imperfective; perfective \emph{\ru{обогнать}} (\emph{\ru{обгоню}, \ru{обгонишь}}). Takes the accusative: \emph{\ru{обгонять} \ru{грузовик}}.

\begin{cousinweb}
\emph{\ru{обгонять}} is built from \emph{\ru{об}-} 'round, past' and \emph{\ru{гонять}} 'to drive'. \emph{\ru{Гонять}} is inherited from Proto-Slavic \emph{ganjati} (EtymDB).
\end{cousinweb}

\begin{grammarlens}[title={What you've built}]
One more word for telling, explaining and describing at B1.
\end{grammarlens}

\subsection*{Your turn}
\begin{itemize}
\raggedright
  \item \emph{Say it:} \emph{obgonyát}
  \item \emph{Say it:} \emph{obgonyát}, once more
  \item \emph{Say it:} a sentence of your own with \emph{obgonyát}
  \item \emph{Recall:} say the Russian for a seat, then the Russian for a windscreen wiper, then say \emph{obgonyát} again
\end{itemize}

\begin{tcolorbox}[breakable,colback=teal!4,colframe=teal!35!black,title={Before you move on}]
What does \ru{обгонять} mean? (\textquotedblleft{}To overtake\textquotedblright{}.) Say it once more.
\end{tcolorbox}
\section[\texorpdfstring{\ru{перестраиваться} (perestráivatsya)}{перестраиваться (perestráivatsya)}]{\ru{перестраиваться} (perestráivatsya) — to change lanes; to adapt}

\label{lesson:RU-C407-perestraivatsya}



\begin{quote}
Before the new one: say the Russian for a windscreen wiper, then the Russian for to overtake.
\end{quote}

\subsection*{You'll want to know: \ru{перестраиваться}}
\textbf{\ru{перестраиваться}} (\emph{perestráivatsya}) — \textquotedblleft{}to change lanes; to adapt\textquotedblright{}.

\begin{quote}
\ru{Перед} \ru{перекрёстком} \ru{она} \ru{перестраивается} \ru{в} \ru{левый} \ru{ряд}.
\end{quote}

\begin{quote}
\emph{Péred perekryóstkom oná perestráivayetsya v lévyy ryad.}
\end{quote}

\begin{quote}
Before the crossroads she moves into the left lane.
\end{quote}

Imperfective; perfective \emph{\ru{перестроиться}}. \emph{\ru{Перестраиваться} \ru{в} \ru{ряд}} — to move into a lane.

\begin{cousinweb}
\emph{\ru{перестраиваться}} is built from \emph{\ru{пере}-} 'across, again', \emph{\ru{строить}} 'to build, to set in rows' and \emph{-\ru{ся}}. EtymDB does not list the verb; its parts are ordinary Russian.
\end{cousinweb}

\begin{grammarlens}[title={What you've built}]
One more word for telling, explaining and describing at B1.
\end{grammarlens}

\subsection*{Your turn}
\begin{itemize}
\raggedright
  \item \emph{Say it:} \emph{perestráivatsya}
  \item \emph{Say it:} \emph{perestráivatsya}, once more
  \item \emph{Say it:} a sentence of your own with \emph{perestráivatsya}
  \item \emph{Recall:} say the Russian for a windscreen wiper, then the Russian for to overtake, then say \emph{perestráivatsya} again
\end{itemize}

\begin{tcolorbox}[breakable,colback=teal!4,colframe=teal!35!black,title={Before you move on}]
What does \ru{перестраиваться} mean? (\textquotedblleft{}To change lanes; to adapt\textquotedblright{}.) Say it once more.
\end{tcolorbox}
\section[\texorpdfstring{\ru{сворачивать} (svoráchivat)}{сворачивать (svoráchivat)}]{\ru{сворачивать} (svoráchivat) — to turn (off a road), to turn (into a street)}

\label{lesson:RU-C407-svorachivat}



\begin{quote}
Before the new one: say the Russian for to overtake, then the Russian for to change lanes.
\end{quote}

\subsection*{You'll want to know: \ru{сворачивать}}
\textbf{\ru{сворачивать}} (\emph{svoráchivat}) — \textquotedblleft{}to turn (off a road), to turn (into a street)\textquotedblright{}.

\begin{quote}
\ru{После} \ru{моста} \ru{мы} \ru{всегда} \ru{сворачиваем} \ru{направо}.
\end{quote}

\begin{quote}
\emph{Pósle mósta my vsegdá svoráchivayem naprávo.}
\end{quote}

\begin{quote}
After the bridge we always turn right.
\end{quote}

Imperfective; perfective \emph{\ru{свернуть}}. \emph{\ru{Свернуть} \ru{налево}} — to turn left.

\begin{cousinweb}
\emph{\ru{сворачивать}} is built from \emph{\ru{с}-} and \emph{-\ru{ворачивать}}, on the root of \emph{\ru{вертеть}} 'to turn'. EtymDB traces \emph{\ru{вертеть}} to Proto-Slavic \emph{vĭrtěti} and the PIE root \emph{wert-} 'to turn'.
\end{cousinweb}

\begin{grammarlens}[title={What you've built}]
One more word for telling, explaining and describing at B1.
\end{grammarlens}

\subsection*{Your turn}
\begin{itemize}
\raggedright
  \item \emph{Say it:} \emph{svoráchivat}
  \item \emph{Say it:} \emph{svoráchivat}, once more
  \item \emph{Say it:} a sentence of your own with \emph{svoráchivat}
  \item \emph{Recall:} say the Russian for to overtake, then the Russian for to change lanes, then say \emph{svoráchivat} again
\end{itemize}

\begin{tcolorbox}[breakable,colback=teal!4,colframe=teal!35!black,title={Before you move on}]
What does \ru{сворачивать} mean? (\textquotedblleft{}To turn (off a road), to turn (into a street)\textquotedblright{}.) Say it once more.
\end{tcolorbox}
\section[\texorpdfstring{\ru{разворачиваться} (razvoráchivatsya)}{разворачиваться (razvoráchivatsya)}]{\ru{разворачиваться} (razvoráchivatsya) — to turn round, to make a U-turn}

\label{lesson:RU-C407-razvorachivatsya}



\begin{quote}
Before the new one: say the Russian for to change lanes, then the Russian for to turn.
\end{quote}

\subsection*{You'll want to know: \ru{разворачиваться}}
\textbf{\ru{разворачиваться}} (\emph{razvoráchivatsya}) — \textquotedblleft{}to turn round, to make a U-turn\textquotedblright{}.

\begin{quote}
\ru{Здесь} \ru{нельзя} \ru{разворачиваться}.
\end{quote}

\begin{quote}
\emph{Zdes nelzyá razvoráchivatsya.}
\end{quote}

\begin{quote}
You are not allowed to turn round here.
\end{quote}

Imperfective; perfective \emph{\ru{развернуться}}. \emph{\ru{Разворот}} — a U-turn.

\begin{cousinweb}
\emph{\ru{разворачиваться}} is built from \emph{\ru{раз}-}, \emph{-\ru{ворачивать}}, on the root of \emph{\ru{вертеть}} 'to turn', and \emph{-\ru{ся}}. EtymDB traces \emph{\ru{вертеть}} to Proto-Slavic \emph{vĭrtěti} and the PIE root \emph{wert-}.
\end{cousinweb}

\begin{grammarlens}[title={What you've built}]
One more word for telling, explaining and describing at B1.
\end{grammarlens}

\subsection*{Your turn}
\begin{itemize}
\raggedright
  \item \emph{Say it:} \emph{razvoráchivatsya}
  \item \emph{Say it:} \emph{razvoráchivatsya}, once more
  \item \emph{Say it:} a sentence of your own with \emph{razvoráchivatsya}
  \item \emph{Recall:} say the Russian for to change lanes, then the Russian for to turn, then say \emph{razvoráchivatsya} again
\end{itemize}

\begin{tcolorbox}[breakable,colback=teal!4,colframe=teal!35!black,title={Before you move on}]
What does \ru{разворачиваться} mean? (\textquotedblleft{}To turn round, to make a U-turn\textquotedblright{}.) Say it once more.
\end{tcolorbox}
\section[\texorpdfstring{\ru{обочина} (obóchina)}{обочина (obóchina)}]{\ru{обочина} (obóchina) — a hard shoulder, a roadside}

\label{lesson:RU-C407-obochina}



\begin{quote}
Before the new one: say the Russian for to turn, then the Russian for to turn round.
\end{quote}

\subsection*{You'll want to know: \ru{обочина}}
\textbf{\ru{обочина}} (\emph{obóchina}) — \textquotedblleft{}a hard shoulder, a roadside\textquotedblright{}.

\begin{quote}
\ru{Сломанная} \ru{машина} \ru{стояла} \ru{на} \ru{обочине}.
\end{quote}

\begin{quote}
\emph{Slómannaya mashína stoyála na obóchine.}
\end{quote}

\begin{quote}
A broken-down car was standing on the hard shoulder.
\end{quote}

Feminine. \emph{\ru{Съехать} \ru{на} \ru{обочину}} — to pull over onto the shoulder.

\begin{cousinweb}
\emph{\ru{обочина}} is usually traced to the old adverb \emph{\ru{обочь}} 'alongside', built on \emph{\ru{бок}} 'a side': the strip at the side of the road. EtymDB does not list the word, so the line is hedged.
\end{cousinweb}

\begin{grammarlens}[title={What you've built}]
One more word for telling, explaining and describing at B1.
\end{grammarlens}

\subsection*{Your turn}
\begin{itemize}
\raggedright
  \item \emph{Say it:} \emph{obóchina}
  \item \emph{Say it:} \emph{obóchina}, once more
  \item \emph{Say it:} a sentence of your own with \emph{obóchina}
  \item \emph{Recall:} say the Russian for to turn, then the Russian for to turn round, then say \emph{obóchina} again
\end{itemize}

\begin{tcolorbox}[breakable,colback=teal!4,colframe=teal!35!black,title={Before you move on}]
What does \ru{обочина} mean? (\textquotedblleft{}A hard shoulder, a roadside\textquotedblright{}.) Say it once more.
\end{tcolorbox}
